\documentclass[11pt,a4paper]{article}

\usepackage[T1]{fontenc}
\usepackage[utf8]{inputenc}
\usepackage[french]{babel}
\usepackage{lmodern}
\usepackage{geometry}
\usepackage{xcolor}
\usepackage{booktabs}
\usepackage{tabularx}
\usepackage{array}
\usepackage{longtable}
\usepackage{listings}
\usepackage{enumitem}
\usepackage{fancyhdr}
\usepackage{hyperref}
\usepackage{microtype}
\usepackage{graphicx}
\usepackage{amsmath}

\geometry{top=2.2cm,bottom=2.3cm,left=2.3cm,right=2.3cm}
\setlength{\parindent}{0pt}
\setlength{\parskip}{0.6em}
\setlength{\headheight}{14pt}
\renewcommand{\arraystretch}{1.22}
\setlist[itemize]{leftmargin=1.5em,itemsep=0.25em,topsep=0.3em}

\definecolor{bleu}{HTML}{17365D}
\definecolor{bleuclair}{HTML}{EAF1F8}
\definecolor{gris}{HTML}{555555}
\definecolor{grisclair}{HTML}{F3F4F6}
\definecolor{vert}{HTML}{246B45}
\definecolor{rouge}{HTML}{9C2F2F}

\hypersetup{
  colorlinks=true,
  linkcolor=bleu,
  urlcolor=bleu,
  pdftitle={Piles generiques validation d expressions et bibliotheque dynamique},
  pdfauthor={Hamri Meriem}
}

\pagestyle{fancy}
\fancyhf{}
\lhead{Piles génériques et bibliothèque dynamique}
\rhead{Hamri Meriem}
\cfoot{\thepage}
\renewcommand{\headrulewidth}{0.4pt}

\lstdefinestyle{cppstyle}{
  language=C++,
  basicstyle=\ttfamily\small,
  keywordstyle=\color{bleu}\bfseries,
  commentstyle=\color{vert},
  stringstyle=\color{rouge},
  numbers=left,
  numberstyle=\tiny\color{gris},
  stepnumber=1,
  numbersep=8pt,
  frame=single,
  rulecolor=\color{gray!35},
  backgroundcolor=\color{grisclair},
  breaklines=true,
  showstringspaces=false,
  tabsize=4,
  columns=fullflexible,
  keepspaces=true
}

\lstdefinestyle{console}{
  basicstyle=\ttfamily\small,
  frame=single,
  rulecolor=\color{gray!35},
  backgroundcolor=\color{black!3},
  breaklines=true,
  showstringspaces=false,
  columns=fullflexible
}

\newcommand{\fichier}[1]{\texttt{#1}}
\newcommand{\complexite}[1]{\ensuremath{\mathcal{O}(#1)}}

\begin{document}

\begin{titlepage}
  \centering
  \vspace*{2.0cm}
  {\Large Rapport de travaux pratiques en C++\par}
  \vspace{1.0cm}
  {\color{bleu}\rule{0.85\textwidth}{1.2pt}\par}
  \vspace{0.8cm}
  {\Huge\bfseries Piles génériques et validation d'expressions\par}
  \vspace{0.35cm}
  {\LARGE Bibliothèque dynamique DLL, complexité et compilation\par}
  \vspace{0.8cm}
  {\color{bleu}\rule{0.85\textwidth}{1.2pt}\par}
  \vfill
  \begin{tabular}{rl}
    \textbf{Réalisé par} & Hamri Meriem \\
    \textbf{Langage} & C++17 \\
    \textbf{Compilateur vérifié} & GCC MinGW-w64 16.2.0 \\
    \textbf{Année universitaire} & 2026--2027
  \end{tabular}
  \vfill
  {\large Version mise à jour du projet\par}
\end{titlepage}

\tableofcontents
\newpage

\section{Introduction}

Ce projet met en œuvre deux piles génériques en C++ et les utilise pour vérifier la syntaxe d'une expression mathématique. La version actuelle ne réalise plus la conversion d'une expression infixe vers une expression préfixée. Elle sépare désormais le traitement en deux composants : une bibliothèque dynamique, \fichier{bibliotheque.dll}, contient les piles et l'algorithme de validation, tandis que \fichier{application.exe} assure la saisie, appelle la DLL et affiche le résultat avec le temps de réponse.

Cette organisation permet d'étudier plusieurs notions complémentaires : la programmation générique avec les \textit{templates}, la structure LIFO d'une pile, l'export de fonctions dans une DLL, l'édition de liens, le chargement dynamique et les différentes transformations réalisées par le compilateur. Le programme annonce une complexité temporelle en \complexite{n}; cette affirmation est justifiée dans le rapport et comparée à des mesures expérimentales.

\section{Objectifs et comportement attendu}

L'application lit une chaîne contenant des nombres, des identificateurs, des opérateurs et des délimiteurs. Elle accepte les opérateurs binaires \texttt{+}, \texttt{-}, \texttt{*}, \texttt{/} et \texttt{\string^}, ainsi que les groupes \texttt{()}, \texttt{[]} et \texttt{\{\}}. Elle doit notamment contrôler :

\begin{itemize}
  \item la correspondance et l'ordre des délimiteurs;
  \item l'alternance cohérente entre opérandes et opérateurs;
  \item la validité des identificateurs, des entiers et des nombres décimaux;
  \item l'absence de caractère interdit, de groupe incomplet ou d'expression vide.
\end{itemize}

Le signe unaire, par exemple \texttt{-x}, n'est pas distingué d'un opérateur binaire. Les appels de fonctions et la multiplication implicite ne sont pas gérés. Il s'agit donc d'un validateur syntaxique ciblé, et non d'un évaluateur mathématique complet.

\section{Organisation du projet}

\begin{table}[ht]
\centering
\caption{Rôle des fichiers de la version actuelle}
\begin{tabularx}{\textwidth}{>{\ttfamily}p{3.4cm} X}
\toprule
\normalfont\textbf{Fichier} & \textbf{Rôle} \\
\midrule
Pile.h & Déclaration des deux classes génériques \texttt{PileTableau<T>} et \texttt{PileListe<T>}. \\
Pile.cpp & Implémentation des piles et instanciations explicites des types placés dans la DLL. \\
Expression.h & Interface publique de la validation : correspondance, opérateurs et \texttt{verifierExpression}. \\
Expression.cpp & Analyse caractère par caractère et utilisation d'une \texttt{PileListe<char>} pour les ouvrants. \\
Export.h & Définition portable de la macro \texttt{BIBLIOTHEQUE\_API}. \\
main.cpp & Saisie, appel de la DLL, chronométrage et affichage de la complexité. \\
build\_dll.bat & Automatisation de la construction de la DLL et de l'exécutable. \\
bibliotheque.dll & Code partagé chargé au moment de l'exécution. \\
libbibliotheque.a & Bibliothèque d'importation utilisée par l'éditeur de liens de MinGW. \\
application.exe & Programme utilisateur dépendant de \texttt{bibliotheque.dll}. \\
\bottomrule
\end{tabularx}
\end{table}

Le fichier \fichier{pile.exe} présent dans le dossier n'est pas produit par \fichier{build\_dll.bat}. Il correspond à un ancien livrable et n'intervient pas dans la chaîne de construction actuelle.

\section{Conception des piles génériques}

\subsection{Interface commune}

Les deux piles sont des classes paramétrées par un type \texttt{T}. Elles exposent la même logique : tester si la pile est vide, empiler une valeur, dépiler la valeur du sommet et consulter le sommet. Une tentative de lecture ou de retrait sur une pile vide déclenche \texttt{std::out\_of\_range}.

\begin{lstlisting}[style=cppstyle,caption={Extrait de l'interface générique}]
template <typename T>
class BIBLIOTHEQUE_API PileListe {
private:
    struct Noeud {
        T valeur;
        Noeud* suivant;
    };
    Noeud* tete;

public:
    PileListe();
    ~PileListe();
    bool estVide() const;
    void empiler(const T& valeur);
    T depiler();
    T sommet() const;
};
\end{lstlisting}

Les constructeurs de copie et opérateurs d'affectation sont supprimés. Cette décision évite une copie superficielle des pointeurs, qui pourrait provoquer des doubles libérations. Une copie correcte demanderait une duplication profonde de tous les éléments ou de tous les maillons.

\subsection{Pile fondée sur un tableau dynamique}

\texttt{PileTableau<T>} alloue initialement deux cases. Lorsque le tableau est plein, \texttt{agrandir()} double la capacité, copie les éléments dans une nouvelle zone, puis libère l'ancienne zone. Les éléments occupent une mémoire contiguë, ce qui favorise généralement la localité de cache.

Une opération \texttt{empiler} ordinaire coûte \complexite{1}. Un agrandissement coûte ponctuellement \complexite{k}, avec $k$ le nombre d'éléments déjà stockés. Comme la capacité double, le coût amorti d'une suite d'empilements reste \complexite{1} par opération. \texttt{depiler}, \texttt{sommet} et \texttt{estVide} sont en \complexite{1}. L'espace utilisé est \complexite{k}.

\subsection{Pile fondée sur une liste chaînée}

\texttt{PileListe<T>} représente chaque élément par un maillon contenant une valeur et un pointeur vers le maillon suivant. La tête est le sommet de la pile. Empiler ajoute une nouvelle tête; dépiler détache et détruit la tête actuelle. Ces deux opérations sont en \complexite{1} dans tous les cas.

Cette représentation n'effectue aucun redimensionnement, mais chaque empilement réalise une allocation dynamique. Elle consomme aussi un pointeur supplémentaire par élément et disperse les maillons en mémoire. Son espace est \complexite{k}. Dans la validation actuelle, cette pile stocke uniquement les délimiteurs ouverts, donc $k$ est au plus égal à la profondeur d'imbrication.

\begin{table}[ht]
\centering
\caption{Comparaison des deux réalisations}
\begin{tabularx}{\textwidth}{l c c X}
\toprule
\textbf{Opération} & \textbf{Tableau} & \textbf{Liste} & \textbf{Commentaire} \\
\midrule
\texttt{empiler} & \complexite{1} amorti & \complexite{1} & Le tableau peut occasionnellement être redimensionné. \\
\texttt{depiler} & \complexite{1} & \complexite{1} & La liste libère un maillon. \\
\texttt{sommet} & \complexite{1} & \complexite{1} & Accès au dernier élément ou à la tête. \\
Mémoire & \complexite{k} & \complexite{k} & La liste ajoute un pointeur et une allocation par valeur. \\
Localité & Bonne & Plus faible & Le tableau est contigu; la liste est dispersée. \\
\bottomrule
\end{tabularx}
\end{table}

\subsection{Templates et instanciations explicites}

Les implémentations des templates se trouvent dans \fichier{Pile.cpp}. Comme le compilateur doit connaître le code d'un template au moment où un type concret est demandé, le fichier termine par des instanciations explicites :

\begin{lstlisting}[style=cppstyle,caption={Types réellement générés dans la DLL}]
template class PileTableau<int>;
template class PileListe<char>;
template class PileListe<std::string>;
\end{lstlisting}

Cette méthode réduit les définitions visibles dans l'en-tête, mais elle impose une contrainte : pour employer un nouveau type, par exemple \texttt{PileTableau<double>}, il faut ajouter son instanciation dans \fichier{Pile.cpp} et reconstruire la DLL. Une autre conception courante consiste à placer toute l'implémentation du template dans le fichier d'en-tête.

\section{Validation d'une expression}

\subsection{Principe de l'algorithme}

La fonction \texttt{verifierExpression} parcourt la chaîne de gauche à droite. La variable booléenne \texttt{attendOperande} joue le rôle d'un petit automate :

\begin{itemize}
  \item si elle vaut \texttt{true}, le prochain élément doit être un opérande ou un délimiteur ouvrant;
  \item si elle vaut \texttt{false}, le prochain élément doit être un opérateur, un délimiteur fermant ou la fin de la chaîne.
\end{itemize}

Un identificateur commence par une lettre ou un souligné, puis peut contenir des caractères alphanumériques ou des soulignés. Un nombre contient une suite de chiffres et éventuellement un point suivi d'au moins un chiffre. Les espaces sont ignorés.

Chaque ouvrant est empilé. Lorsqu'un fermant apparaît, il doit correspondre au sommet de la pile. Ce fonctionnement LIFO est indispensable : le dernier groupe ouvert doit être le premier refermé. Par exemple, dans \texttt{\{(x+y)+z)*5}, le second \texttt{)} rencontre \texttt{\{} au sommet; le programme signale alors des délimiteurs mal ordonnés.

\subsection{Extrait commenté}

\begin{lstlisting}[style=cppstyle,caption={Contrôle des délimiteurs fermants}]
else if (caractere == ')' || caractere == ']' || caractere == '}') {
    // Un fermant ne peut pas remplacer un operande attendu.
    if (attendOperande) {
        erreur = "Operande manquant avant le fermant";
        return false;
    }

    // Le dernier ouvrant doit etre de la meme famille.
    if (ouvrants.estVide() ||
        !correspondent(ouvrants.sommet(), caractere)) {
        erreur = "Parentheses ou delimiteurs mal ordonnes";
        return false;
    }

    ouvrants.depiler();
    attendOperande = false;
    ++i;
}
\end{lstlisting}

Le retour anticipé est utile : dès qu'une erreur est certaine, le reste de la chaîne n'a plus besoin d'être parcouru. Dans le pire cas, toutefois, l'expression entière est lue.

\section{Complexité de la validation}

Soit $n$ le nombre de caractères de l'expression et $d$ la profondeur maximale d'imbrication des délimiteurs.

\begin{itemize}
  \item Chaque position est visitée au plus un nombre constant de fois. Les boucles internes qui lisent un identificateur ou un nombre font progresser le même indice \texttt{i}; elles ne recommencent pas au début. Le temps total est donc \complexite{n} dans le pire cas.
  \item En cas d'erreur précoce, le temps peut être inférieur à $n$, mais la borne asymptotique du pire cas reste \complexite{n}.
  \item La pile contient uniquement les ouvrants non encore refermés. Sa mémoire auxiliaire est \complexite{d}; puisque $d \leq n$, la borne générale est \complexite{n}.
  \item Les chaînes de message d'erreur ont une taille bornée par rapport à l'entrée et ne modifient pas l'ordre de complexité.
\end{itemize}

L'affichage de \texttt{Complexite temporelle (pire cas) : O(n)} dans \fichier{main.cpp} est donc cohérent avec l'algorithme. La valeur $n$ affichée correspond à \texttt{expression.size()}, c'est-à-dire au nombre d'octets de la chaîne. Pour les caractères ASCII acceptés par ce programme, ce nombre correspond aussi au nombre de caractères visibles.

\section{Mesure du temps d'exécution}

\subsection{Méthode utilisée dans le programme}

Le chronométrage emploie \texttt{std::chrono::steady\_clock}, une horloge monotone adaptée aux durées. Il commence après la saisie et s'arrête immédiatement après l'appel à \texttt{verifierExpression}. Le temps mesuré exclut donc l'attente de l'utilisateur, l'affichage, et la majeure partie du coût de démarrage de l'exécutable.

\begin{lstlisting}[style=cppstyle,caption={Zone effectivement chronométrée}]
const auto debut = std::chrono::steady_clock::now();
const bool estValide = verifierExpression(expression, erreur);
const auto fin = std::chrono::steady_clock::now();

const double dureeMicrosecondes =
    std::chrono::duration<double, std::micro>(fin - debut).count();
\end{lstlisting}

Une mesure unique sur une chaîne très courte peut être influencée par la précision de l'horloge, l'ordonnancement de Windows, les caches et l'allocation mémoire. Pour obtenir une lecture plus stable, quatre tailles valides ont été exécutées 30 fois chacune. Chaque expression est de la forme \texttt{a+a+...+a}. Les valeurs suivantes proviennent de l'exécutable actuel, compilé avec GCC MinGW-w64 16.2.0.

\begin{table}[ht]
\centering
\caption{Temps observés pour l'appel à la fonction de validation}
\begin{tabular}{r r r r r}
\toprule
\textbf{Taille $n$} & \textbf{Minimum (µs)} & \textbf{Médiane (µs)} & \textbf{Moyenne (µs)} & \textbf{Maximum (µs)} \\
\midrule
9    & 4,9   & 5,4   & 6,05  & 18,2 \\
99   & 7,8   & 8,6   & 9,28  & 23,4 \\
999  & 38,6  & 39,7  & 47,31 & 65,3 \\
9 999 & 345,6 & 349,7 & 365,00 & 542,2 \\
\bottomrule
\end{tabular}
\end{table}

La médiane passe de 39,7 µs pour 999 caractères à 349,7 µs pour 9 999 caractères. La taille est multipliée par environ dix et le temps par environ 8,8, ce qui est compatible avec une croissance linéaire. Les petites entrées ne suivent pas une proportion parfaite, car le coût fixe de l'appel, de l'horloge et des allocations devient comparable au travail utile.

Ces chiffres décrivent une machine et une exécution données; ils ne constituent pas une garantie universelle. Une campagne plus rigoureuse placerait plusieurs milliers d'appels dans un même processus, séparerait les phases d'échauffement et compilerait aussi avec optimisation, par exemple \texttt{-O2}.

\section{Bibliothèque dynamique et exécutable}

\subsection{Exportation et importation}

La macro de \fichier{Export.h} change selon le composant compilé. Pendant la construction de la DLL, \texttt{BIBLIOTHEQUE\_EXPORTS} est défini et les symboles reçoivent \texttt{\_\_declspec(dllexport)}. Pendant la compilation de l'application, la macro n'est pas définie et les mêmes déclarations deviennent \texttt{\_\_declspec(dllimport)}.

\begin{lstlisting}[style=cppstyle,caption={Macro de visibilité sous Windows}]
#if defined(_WIN32) || defined(__CYGWIN__)
    #ifdef BIBLIOTHEQUE_EXPORTS
        #define BIBLIOTHEQUE_API __declspec(dllexport)
    #else
        #define BIBLIOTHEQUE_API __declspec(dllimport)
    #endif
#else
    #define BIBLIOTHEQUE_API
#endif
\end{lstlisting}

L'inspection de la DLL confirme la présence des symboles \texttt{correspondent}, \texttt{estOperateur}, \texttt{verifierExpression} et des instanciations explicites des piles. L'exécutable déclare directement une dépendance envers \fichier{bibliotheque.dll}.

\subsection{Rôle des trois livrables}

\begin{description}[style=nextline,leftmargin=0pt]
  \item[\fichier{bibliotheque.dll}] contient le code machine partagé. Il doit être accessible au chargement, ici dans le même dossier que l'exécutable.
  \item[\fichier{libbibliotheque.a}] est une bibliothèque d'importation MinGW. Au moment de l'édition de liens de l'application, elle décrit les symboles qui seront fournis par la DLL. Elle ne remplace pas la DLL à l'exécution.
  \item[\fichier{application.exe}] contient \texttt{main} et une table d'importation. Le chargeur de Windows résout cette table vers \fichier{bibliotheque.dll} avant l'exécution du programme.
\end{description}

Si la DLL est absente ou incompatible, l'application ne peut pas démarrer correctement. Les DLL du runtime MinGW, notamment \fichier{libstdc++-6.dll} et \fichier{libgcc\_s\_seh-1.dll}, doivent également être trouvées dans l'environnement ou à proximité du programme.

\section{Processus complet de compilation}

\subsection{Vue d'ensemble des transformations}

La commande unique de compilation masque plusieurs étapes. Pour un fichier C++, la chaîne peut être représentée ainsi :

\begin{center}
\fichier{.cpp + .h}
$\longrightarrow$ préprocesseur
$\longrightarrow$ \fichier{.ii}
$\longrightarrow$ compilateur
$\longrightarrow$ \fichier{.s}
$\longrightarrow$ assembleur
$\longrightarrow$ \fichier{.o}
$\longrightarrow$ éditeur de liens
$\longrightarrow$ \fichier{.dll / .a / .exe}
\end{center}

L'extension \fichier{.i} est traditionnellement utilisée pour un fichier C prétraité, tandis que \fichier{.ii} indique du C++ prétraité. Le contenu joue un rôle similaire, mais l'extension aide GCC à choisir le langage. Comme le projet est en C++, \fichier{.ii} est l'extension correcte.

\subsection{Étape 1 prétraitement vers .ii}

L'option \texttt{-E} exécute seulement le préprocesseur. Les \texttt{\#include} sont développés, les macros sont remplacées et les branches conditionnelles sont choisies. Pour la DLL, il faut déjà définir \texttt{BIBLIOTHEQUE\_EXPORTS} afin de sélectionner \texttt{dllexport}.

\begin{lstlisting}[style=console,caption={Prétraitement C++}]
g++ -std=c++17 -DBIBLIOTHEQUE_EXPORTS -E Pile.cpp       -o Pile.ii
g++ -std=c++17 -DBIBLIOTHEQUE_EXPORTS -E Expression.cpp -o Expression.ii
g++ -std=c++17 -E main.cpp -o main.ii
\end{lstlisting}

Les fichiers deviennent volumineux parce qu'ils contiennent aussi les en-têtes de la bibliothèque standard. Lors de la vérification, les tailles obtenues étaient de 793 337 octets pour \fichier{Pile.ii}, 777 125 pour \fichier{Expression.ii} et 1 332 914 pour \fichier{main.ii}.

\subsection{Étape 2 compilation vers .s}

L'option \texttt{-S} traduit le C++ prétraité en langage assembleur. C'est à cette étape que les vérifications de types, la génération des instanciations de templates et les transformations d'optimisation sont réalisées.

\begin{lstlisting}[style=console,caption={Génération de l'assembleur}]
g++ -std=c++17 -S Pile.ii       -o Pile.s
g++ -std=c++17 -S Expression.ii -o Expression.s
g++ -std=c++17 -S main.ii       -o main.s
\end{lstlisting}

Les fichiers assembleur vérifiés avaient respectivement une taille de 61 706, 87 966 et 47 664 octets. Ils restent lisibles comme du texte, mais décrivent des instructions et des symboles propres à l'architecture x86-64.

\subsection{Étape 3 assemblage vers .o}

L'option \texttt{-c} demande la production d'un fichier objet sans lancer l'édition de liens. L'assembleur convertit les instructions textuelles en code machine. Le fichier \fichier{.o} contient du code binaire, des symboles et des informations de relocation, car les adresses définitives ne sont pas encore toutes connues.

\begin{lstlisting}[style=console,caption={Production des objets}]
g++ -c Pile.s       -o Pile.o
g++ -c Expression.s -o Expression.o
g++ -c main.s       -o main.o
\end{lstlisting}

Les tailles mesurées étaient de 50 699 octets pour \fichier{Pile.o}, 59 314 pour \fichier{Expression.o} et 44 632 pour \fichier{main.o}.

\subsection{Étape 4 édition de liens de la DLL}

L'option \texttt{-shared} demande une bibliothèque dynamique. L'éditeur de liens regroupe \fichier{Pile.o} et \fichier{Expression.o}, résout leurs références et crée en même temps la bibliothèque d'importation grâce à \texttt{--out-implib}.

\begin{lstlisting}[style=console,caption={Création de la DLL et de sa bibliothèque d'importation}]
g++ -shared Pile.o Expression.o \
    -Wl,--out-implib,libbibliotheque.a \
    -o bibliotheque.dll
\end{lstlisting}

Les livrables vérifiés mesurent 82 322 octets pour \fichier{bibliotheque.dll} et 22 816 octets pour \fichier{libbibliotheque.a}.

\subsection{Étape 5 édition de liens de l'application}

L'option \texttt{-L.} ajoute le dossier courant aux emplacements recherchés. L'option \texttt{-lbibliotheque} demande un fichier dont le nom correspond à \fichier{libbibliotheque.a}. L'éditeur de liens inscrit alors dans \fichier{application.exe} les importations à résoudre dans la DLL.

\begin{lstlisting}[style=console,caption={Création de l'exécutable}]
g++ main.o -L. -lbibliotheque -o application.exe
\end{lstlisting}

L'exécutable obtenu mesure 68 734 octets. L'inspection de sa table d'importation mentionne explicitement \fichier{bibliotheque.dll}, ce qui confirme que la liaison est dynamique.

\subsection{Étape 6 chargement à l'exécution}

Lorsque l'utilisateur lance \fichier{application.exe}, le chargeur de Windows recherche les DLL nécessaires, les place dans l'espace mémoire du processus puis associe les symboles importés aux adresses exportées. L'appel C++ à \texttt{verifierExpression} peut ensuite atteindre le code de \fichier{bibliotheque.dll}. Cette dernière étape n'est pas une compilation : elle se produit à chaque démarrage du programme.

\subsection{Automatisation par le script batch}

Le fichier \fichier{build\_dll.bat} regroupe les étapes en deux commandes de haut niveau. GCC exécute automatiquement le prétraitement, la compilation et l'assemblage; les fichiers \fichier{.ii}, \fichier{.s} et \fichier{.o} temporaires sont normalement supprimés. Le script fixe également \texttt{TEMP} et \texttt{TMP} au dossier du projet et ajoute l'environnement MinGW-w64 UCRT au \texttt{PATH} lorsqu'il est installé dans \fichier{C:\textbackslash msys64\textbackslash ucrt64\textbackslash bin}.

\begin{lstlisting}[style=console,caption={Utilisation normale du projet}]
build_dll.bat
build_dll.bat run
\end{lstlisting}

La première commande compile uniquement. La seconde compile puis lance l'application. Les options \texttt{-Wall -Wextra -pedantic} activent des diagnostics utiles, tandis que \texttt{-std=c++17} fixe la version du langage.

\section{Résultats fonctionnels}

\begin{table}[ht]
\centering
\caption{Cas exécutés avec l'application actuelle}
\begin{tabularx}{\textwidth}{>{\ttfamily}p{3.7cm} l X}
\toprule
\normalfont\textbf{Expression} & \textbf{Résultat} & \textbf{Interprétation} \\
\midrule
a+b*c & Valide & L'alternance opérande-opérateur est cohérente. La priorité n'est pas évaluée car le programme valide seulement la syntaxe. \\
\{(x+y)+z)*5 & Invalide & Le second \texttt{)} ne correspond pas à l'accolade située au sommet de la pile. \\
\{(x+y)+z\}*5 & Valide & Chaque groupe est fermé dans l'ordre inverse de son ouverture. \\
a+*b & Invalide & Après \texttt{+}, un opérande était attendu; \texttt{*} est donc rejeté. \\
12.5*(x\_1-3) & Valide & Le décimal et l'identificateur contenant un souligné sont reconnus. \\
Chaîne vide & Invalide & Aucun élément syntaxique n'a été rencontré. \\
\bottomrule
\end{tabularx}
\end{table}

Les messages sont suffisamment précis pour identifier la catégorie de l'erreur. Certains pourraient encore être améliorés en ajoutant systématiquement la position fautive, comme cela est déjà fait pour un opérateur manquant ou un caractère interdit.

\section{Analyse critique et améliorations possibles}

La séparation DLL/application rend le projet modulaire et illustre correctement le mécanisme d'importation. L'algorithme de validation est simple, lisible et linéaire. La pile chaînée convient à la profondeur d'imbrication, même si une petite pile fondée sur un tableau pourrait réduire le nombre d'allocations.

Plusieurs améliorations seraient possibles :

\begin{itemize}
  \item conserver dans la pile la position de chaque ouvrant afin de signaler précisément un groupe non fermé;
  \item ajouter une taille à \texttt{PileListe} et harmoniser davantage l'interface des deux piles;
  \item employer \texttt{std::unique\_ptr} pour rendre la propriété des maillons explicite et plus sûre;
  \item ajouter les constructeurs de déplacement, ou placer les implémentations de templates dans l'en-tête pour faciliter l'emploi de nouveaux types;
  \item créer une suite de tests automatisés couvrant les limites, les nombres mal formés et les imbrications profondes;
  \item réaliser un benchmark dans un seul processus, avec de nombreuses répétitions et une version optimisée \texttt{-O2}.
\end{itemize}

\section{Conclusion}

La version actuelle du projet ne se limite plus à démontrer deux piles génériques. Elle montre comment isoler un traitement réutilisable dans une bibliothèque dynamique, comment exporter son interface et comment lier une application cliente. La validation parcourt chaque caractère une seule fois au sens asymptotique : sa complexité temporelle est \complexite{n} et sa mémoire auxiliaire est \complexite{d}, où $d$ est la profondeur d'imbrication.

Les mesures réalisées sur des expressions de 9 à 9 999 caractères suivent la tendance linéaire attendue. La reconstruction manuelle des étapes \fichier{.ii}, \fichier{.s}, \fichier{.o}, \fichier{.dll}, \fichier{.a} et \fichier{.exe} confirme aussi le rôle de chaque outil : préprocesseur, compilateur, assembleur, éditeur de liens et chargeur dynamique. Le projet fournit ainsi un exemple complet reliant structures de données, analyse syntaxique, complexité et chaîne de compilation C++ sous Windows.

\appendix
\section{Commandes de vérification utiles}

Les dépendances et symboles peuvent être inspectés avec les outils MinGW :

\begin{lstlisting}[style=console]
objdump -p application.exe
objdump -p bibliotheque.dll
nm -C -g bibliotheque.dll
\end{lstlisting}

\texttt{objdump -p} affiche notamment la table d'importation. \texttt{nm -C -g} affiche les symboles globaux et démangle les noms C++, ce qui rend visibles les types instanciés et les fonctions exportées.

\end{document}
